\answer{ex:21:1}{(a)~$T=\tau_{\text{rec}}\ln\frac{1}{1-x^*}$: $1.84$~s et $3.68$~s. (b)~De $3.36$ à $4.24$~s: la sensibilité $\dd T/\dd x^*=\tau_{\text{rec}}/(1-x^*)=80$~s; quand $x^*\to1$, les salves deviennent irrégulières.}
\answer{ex:21:2}{(a)~$8.6$~W, comme dans~\eqref{eq:energy}. (b)~$205$~Mbit/s, soit $0.2$~W: $2\cdot10^3$ fois plus que la culture ($10^{-4}$~W).}
\answer{ex:21:3}{$x_{\text{st}}=1/(1+Ur_b\tau_{\text{rec}})<x^*$ pour $r_b>(1/x^*-1)/(U\tau_{\text{rec}})=0.28$~Hz si $x^*=0.9$ et $0.025$~Hz si $x^*=0.99$; la force de la synapse tombe à $x_{\text{st}}$, c'est-à-dire de $10\%$ et de $1\%$.}
\answer{ex:21:4}{Les $u(t-k)$ sont orthonormés, $m_k=\|P_Vu(t-k)\|^2$, où $P_V$ est le projecteur sur l'enveloppe linéaire des sorties, de dimension $N$; par l'inégalité de Bessel, $\sum_k m_k\le N$.}
\answer{ex:21:5}{$\mathrm{MC}\approx9$--$11$ avec $\tanh$ et $15$--$18$ pour le réservoir linéaire (trois graines), tous deux $\le30$: la non-linéarité se paie en mémoire.}
\answer{ex:21:6}{(a)~$n\le102$. (b)~De $5.6$.}
\answer{ex:21:7}{Problème ouvert. Contrôle clé: geler la lecture et vérifier si l'amélioration subsiste après une pause sans stimulation; sinon, c'est l'état qui a changé, et non les poids.}
